\chapter{包、临时环境与开发环境}
\chapterlead{传统发行版在系统层面装包，NixOS 可以按需构造一次性环境。本章说明四种装包方式的差别、如何为单个项目固定依赖，以及如何把开发环境写成可提交的文件。}

\section{起步：四种方式}

\begin{shellsession}
$ nix shell nixpkgs#jq          # 临时进入一个含 jq 的 shell
$ nix run nixpkgs#cowsay -- hi  # 直接运行，不进入 shell
$ nix profile add nixpkgs#ripgrep   # 装进用户 profile，长期可用
$ nix search nixpkgs "json"     # 搜索包
\end{shellsession}

四种方式的差别在于生命周期与作用范围。

\begin{booktable}{\textwidth}{J{1.25} J{0.75} J{1.0}}
\toprule
方式 & 生效范围 & 适用场景 \\
\midrule
\texttt{environment.\allowbreak{}systemPackages} & 全系统 & 常用工具，希望所有用户可用 \\
\texttt{nix profile add} & 当前用户 & 个人工具，不想改系统配置 \\
\texttt{nix shell} & 当前 shell & 试一个命令，或临时需要某个工具 \\
\texttt{nix run} & 一次调用 & 只运行一次，不需要留在环境里 \\
\bottomrule
\end{booktable}

\begin{notebox}[systemPackages 与 profile 的取舍]
\texttt{environment.\allowbreak{}systemPackages} 写在配置里，随重建生效，可回滚，适合大多数情况。\texttt{nix profile} 装在用户目录下，不进入系统配置，重装系统后需要重新添加。日常工具优先写配置；只在试某个软件时用 \texttt{nix shell}。
\end{notebox}

\section{为项目固定依赖}

项目需要的不只是可执行文件，还包括特定版本的编译器与库。Nix 用一份 \texttt{shell.nix} 或 \texttt{flake.nix} 描述这些依赖。

\begin{nixcode}
# shell.nix
{ pkgs ? import <nixpkgs> {} }:

pkgs.mkShell {
  packages = with pkgs; [
    python312
    python312Packages.requests
    ruff
    mypy
  ];

  shellHook = ''
    echo "Python 环境就绪：$(python --version)"
  '';
}
\end{nixcode}

进入环境：

\begin{shellsession}
$ nix-shell            # 读取 shell.nix
$ nix develop          # 读取 flake.nix 中的 devShell
$ exit                 # 离开，环境消失
\end{shellsession}

环境里只有声明的包，系统的 \texttt{PATH} 被替换。这解决了一个长期问题：项目 A 要 Node 20，项目 B 要 Node 22，两者的依赖不冲突，也不需要版本管理工具。

\section{可复现的版本固定}

\texttt{shell.nix} 里的 \texttt{<nixpkgs>} 指向当前系统的通道，换一台机器可能得到不同版本。要真正可复现，把 nixpkgs 固定到某个提交。

\begin{nixcode}
# flake.nix
{
  inputs.nixpkgs.url = "github:NixOS/nixpkgs/nixos-26.11";

  outputs = { self, nixpkgs }:
    let
      system = "x86_64-linux";
      pkgs = import nixpkgs { inherit system; };
    in {
      devShells.${system}.default = pkgs.mkShell {
        packages = with pkgs; [ python312 ruff ];
      };
    };
}
\end{nixcode}

\begin{shellsession}
$ nix develop
$ nix flake lock        # 生成 flake.lock，记录精确版本
\end{shellsession}

\texttt{flake.lock} 应当提交进仓库。它记录了每个输入的哈希，任何人克隆后执行 \texttt{nix develop} 都得到同一套依赖。这是 Nix 相对其他包管理器最实在的收益。

\section{打包一个没有收录的软件}

nixpkgs 收录了十万量级的包，仍会有漏网者。最小的打包写法是一个 \texttt{derivation}。

\begin{nixcode}
{ pkgs ? import <nixpkgs> {} }:

pkgs.stdenv.mkDerivation {
  pname = "hello-script";
  version = "0.1.0";

  src = ./src;

  installPhase = ''
    mkdir -p $out/bin
    cp hello.sh $out/bin/hello
    chmod +x $out/bin/hello
  '';
}
\end{nixcode}

\texttt{\$out} 是构建产物要落到的 store 路径。构建用 \texttt{nix-build}，结果出现在 \texttt{result} 符号链接下。真实项目的打包需要处理依赖、编译参数、测试，nixpkgs 手册有完整说明；上面的骨架足以把脚本类工具纳入统一管理。

\section{练习}

\begin{enumerate}
  \item 分别用 \texttt{nix shell} 与 \texttt{nix run} 使用 \texttt{jq}，说明两者退出后各自的残留。
  \item 写一个 \texttt{shell.nix}，提供 \texttt{python312} 与 \texttt{ruff}，进入后运行 \texttt{ruff --version}。
  \item 把上题改写成 \texttt{flake.nix} 的 devShell，生成 \texttt{flake.lock}，记录锁定到的 nixpkgs 提交。
  \item 用 \texttt{nix-build} 打包一个 shell 脚本，验证 \texttt{result/bin} 下的可执行文件。
  \item 用 \texttt{nix path-info -Sh} 比较空环境与开发环境的大小差异，解释差额来源。
\end{enumerate}

\begin{checkpoint}
本章结束时，应能：
\begin{itemize}
  \item 说明四种装包方式的生命周期与适用场景；
  \item 用 \texttt{shell.nix} 为项目提供隔离的开发环境；
  \item 用 flake 固定依赖版本，理解 \texttt{flake.lock} 的作用；
  \item 写出一个最小的 derivation，把未收录的脚本纳入 Nix 管理；
  \item 判断什么时候用系统包，什么时候用项目环境。
\end{itemize}
\end{checkpoint}
